\section{Experiments}

\subsection{Setup and correspondence to the theory}

All experiments are synthetic and static. They test the controlled claims one
at a time. We first test whether well-spread prototype cells balance uniform
spherical features. We then separate fixed prototype-cell imbalance from
finite-token sampling noise, test expert spectral prototypes as candidate
centers, test the covariance-isotropy counterexample, and test spectral
flattening only in the elliptical toy family where Proposition~\ref{prop:whitening}
applies. Finally, we include capacity assignment as a forced-count control.

The main metric is LoadCV, the coefficient of variation of empirical expert
usage:
\begin{equation}
 \mathrm{LoadCV}=\frac{\mathrm{std}(\widehat p_1,\ldots,\widehat p_E)}
 {\mathrm{mean}(\widehat p_1,\ldots,\widehat p_E)}.
\end{equation}
For spectral geometry, we report covariance condition number and effective
rank. For spectral expert prototypes, we report
\begin{equation}
 \lambda_e=\frac{\|r_e^\top W_e\|_2}{\|r_e\|_2\|W_e\|_2}.
\end{equation}

\subsection{Prototype geometry controls uniform-sphere usage}

We sample $N=100000$ examples from $\Sphere^{63}$ and route them to $E=8$
experts. Symmetric prototypes are regular-simplex rows, random prototypes are
independent unit vectors, and clustered prototypes are perturbations of one
shared direction.

\begin{table}[t]
\centering
\caption{Uniform-sphere routing with $d=64$ and $E=8$.}
\label{tab:uniform_sphere}
\small
\resizebox{\linewidth}{!}{%
\begin{tabular}{lccc}
\toprule
Prototype geometry & LoadCV $\downarrow$ & Max/min $\downarrow$
& Max abs. error $\downarrow$ \\
\midrule
Symmetric top-1 & 0.0067 & 1.021 & 0.00145 \\
Random top-1 & 0.0657 & 1.227 & 0.01451 \\
Clustered top-1 & 0.1349 & 1.574 & 0.03580 \\
\bottomrule
\end{tabular}
}
\end{table}

\begin{figure}[t]
\centering
\includegraphics[width=0.95\linewidth]{figures/prototype_geometry_load.pdf}
\caption{Uniform spherical features do not alone determine expert usage.
Prototype geometry controls load variation.}
\label{fig:prototype_geometry}
\end{figure}

Table~\ref{tab:uniform_sphere} and Figure~\ref{fig:prototype_geometry} support
the sufficient-condition story: well-spread cells give near-uniform usage,
whereas clustered cells do not.

\subsection{Random prototype imbalance is a geometry floor}

We estimate high-dimensional random-prototype cell masses by Monte Carlo
reference sampling. For a fixed random prototype realization, we compare
geometry error, sampling error, and total usage error as routed batch size
increases.

\begin{table}[t]
\centering
\caption{Random-prototype floor for $d=64$.}
\label{tab:geometry_floor}
\small
\resizebox{\linewidth}{!}{%
\begin{tabular}{ccccc}
\toprule
$E$ & $N$ & Geometry & Sampling & Total \\
\midrule
4 & $10^3$ & 0.093 & 0.074 & 0.096 \\
4 & $5{\cdot}10^4$ & 0.093 & 0.014 & 0.097 \\
8 & $10^3$ & 0.093 & 0.160 & 0.182 \\
8 & $5{\cdot}10^4$ & 0.093 & 0.031 & 0.109 \\
16 & $10^3$ & 0.185 & 0.237 & 0.234 \\
16 & $5{\cdot}10^4$ & 0.185 & 0.044 & 0.194 \\
\bottomrule
\end{tabular}
}
\end{table}

\begin{figure}[t]
\centering
\includegraphics[width=0.98\linewidth]{figures/geometry_floor_sampling.pdf}
\caption{Increasing routed examples reduces sampling noise but does not remove
the single-realization random-prototype geometry floor.}
\label{fig:geometry_floor}
\end{figure}

Table~\ref{tab:geometry_floor} and Figure~\ref{fig:geometry_floor} match the
finite-token decomposition: more examples reduce $\widehat p_e-p_e$, not
$p_e-1/E$.

\subsection{Spectral prototypes are centers, not balance}

We construct expert matrices with prescribed input-side signal directions, not
exact finite-noise singular vectors. For input dimension $d$ and output width
$m$, choose unit vectors $v_e\in\mathbb{R}^d$ and $a_e\in\mathbb{R}^m$, and set
\[
W_e=\sigma_1 v_ea_e^\top+\sigma_0 Z_e,\qquad
Z_e\in\mathbb{R}^{d\times m},\quad (Z_e)_{ij}\sim\mathcal{N}(0,1).
\]
When $\sigma_0=0$, $v_e$ is the top left singular vector; with finite Gaussian
noise it is a perturbative target, so we report measured alignment instead of
assuming exact recovery. The prescribed directions $v_e$ are symmetric, random,
or clustered. We compare exact singular value decomposition (SVD), one-step
MPI, five-step MPI, and random rows.

\begin{table}[t]
\centering
\caption{Spectral alignment recovers expert directions, but usage balance still
depends on prototype spread.}
\label{tab:spectral_prototypes}
\small
\setlength{\tabcolsep}{3pt}
\resizebox{\linewidth}{!}{%
\begin{tabular}{llccc}
\toprule
Geometry & Router & Align. $\uparrow$ & LoadCV $\downarrow$ & Max/min $\downarrow$ \\
\midrule
Sym. & Exact SVD & 1.000 & 0.009 & 1.030 \\
Sym. & MPI-1 & 0.975 & 0.025 & 1.094 \\
Clust. & Exact SVD & 1.000 & 0.161 & 1.631 \\
Clust. & MPI-1 & 0.977 & 0.195 & 1.788 \\
Random rows & None & 0.098 & 0.074 & 1.281 \\
\bottomrule
\end{tabular}
}
\end{table}

\begin{figure}[t]
\centering
\includegraphics[width=0.98\linewidth]{figures/spectral_alignment_not_balance.pdf}
\caption{High spectral alignment does not imply balanced expert usage.}
\label{fig:spectral_alignment}
\end{figure}

Table~\ref{tab:spectral_prototypes} and Figure~\ref{fig:spectral_alignment}
show that expert singular directions can be meaningful centers, but a
center-construction rule is not a load-balance theorem.

\subsection{Isotropic covariance is not enough}

This experiment mirrors Proposition~\ref{prop:isotropic_not_directional}. The
token distribution has equal mass on the four axis directions in $\R^2$, so it
has mean zero and covariance $I/2$. We route these points with three
regular-simplex prototypes at angles $20^\circ$, $140^\circ$, and $260^\circ$.

\begin{table}[t]
\centering
\caption{Covariance isotropy does not imply uniform routing.}
\label{tab:covariance_counterexample}
\small
\resizebox{\linewidth}{!}{%
\begin{tabular}{lcccc}
\toprule
Distribution & Cond. & Eff. rank & Load vector & LoadCV \\
\midrule
Four axis points & 1.000 & 2.000 & $(.25,.50,.25)$ & 0.3536 \\
\bottomrule
\end{tabular}
}
\end{table}

\begin{figure}[t]
\centering
\includegraphics[width=0.95\linewidth]{figures/covariance_counterexample.pdf}
\caption{A flat covariance spectrum does not certify uniform routing.}
\label{fig:covariance_counterexample}
\end{figure}

Table~\ref{tab:covariance_counterexample} and
Figure~\ref{fig:covariance_counterexample} provide the operational version of
the counterexample: condition number 1 coexists with nonuniform load.

\subsection{Whitening repairs an elliptical toy}

We test Proposition~\ref{prop:whitening} in a controlled family. Starting from
$z\sim\mathrm{Unif}(\Sphere^{63})$, we form anisotropic features $h=Dz$ with
$D=\mathrm{diag}(\alpha,1,\ldots,1)$. We compare raw features,
centered-normalized features, partial spectral shrinkage with $\beta=0.5$, and
whitening.

\begin{table}[t]
\centering
\caption{Representation-geometric repair for $\alpha=20$.}
\label{tab:spectrum_uniformization}
\small
\setlength{\tabcolsep}{3pt}
\resizebox{\linewidth}{!}{%
\begin{tabular}{lccc}
\toprule
Representation & Cond. $\downarrow$ & Eff. rank $\uparrow$ & LoadCV $\downarrow$ \\
\midrule
Raw & 419.52 & 2.61 & 1.053 \\
Centered + norm. & 114.13 & 8.95 & 1.053 \\
Partial shrinkage & 14.86 & 47.02 & 0.781 \\
Whitened + norm. & 1.00 & 64.00 & 0.005 \\
\bottomrule
\end{tabular}
}
\end{table}

\begin{figure}[t]
\centering
\includegraphics[width=0.98\linewidth]{figures/spectrum_uniformization.pdf}
\caption{In the elliptical toy setting, spectral flattening moves the
representation back toward the spherical routing model.}
\label{fig:spectrum_uniformization}
\end{figure}

Table~\ref{tab:spectrum_uniformization} and
Figure~\ref{fig:spectrum_uniformization} show that whitening is successful here
because the data are generated by an elliptical distortion of a spherical
source. This is not a claim that arbitrary hidden states are elliptical or
naturally whitened.

\subsection{Capacity assignment is only a control}

With $N=2048,d=64,E=8$, both greedy capacity-feasible assignment and exact
balanced assignment over expanded expert slots enforce zero LoadCV.

\begin{table}[t]
\centering
\caption{Capacity assignment as a non-geometric control.}
\label{tab:balanced_assignment}
\small
\setlength{\tabcolsep}{3pt}
\resizebox{\linewidth}{!}{%
\begin{tabular}{lccc}
\toprule
Geometry & Top-1 LoadCV & Greedy drop & Exact drop \\
\midrule
Symmetric & 0.054 & 0.00091 & 0.00012 \\
Random & 0.122 & 0.00416 & 0.00056 \\
Clustered & 0.128 & 0.00379 & 0.00047 \\
\bottomrule
\end{tabular}
}
\end{table}

\begin{figure}[t]
\centering
\includegraphics[width=0.98\linewidth]{figures/balanced_assignment_cost.pdf}
\caption{Balanced assignment exposes the cost of forcing the count metric. It is
a control for count equality, not the representation mechanism studied here.}
\label{fig:balanced_assignment}
\end{figure}

Table~\ref{tab:balanced_assignment} and Figure~\ref{fig:balanced_assignment}
show why this control belongs in the paper: uniform expert usage can be forced,
but forced usage is not evidence for uniform feature geometry.
